\clearpage
\section{Data Formats}\label{page:data-formats}

\subsection{Integer Data Sizes}
Integer operands use size suffixes to select the low-order subfield of an Rn register and the number of bytes
transferred by memory operands.

Byte, word, and long operations use the low-order subfield of the selected register. Q-sized operations use the full
64-bit register.

Every Rn destination write defines all 64 bits. A Q-sized write uses the complete result. For a B-, W-, or L-sized
write, \texttt{ABS}, \texttt{DIVS}, \texttt{MODS}, \texttt{DIVMODS}, \texttt{MINS}, \texttt{MAXS}, \texttt{SAR},
and the explicit \texttt{EXTS} operations sign-extend the selected-width result to 64 bits. Every other integer
operation, including \texttt{MOV}, neutral modular arithmetic, logical and bit operations, count and Boolean results,
address results, atomic return values, and state-register reads, zero-extends the selected-width result to 64 bits.
This result extension does not change FLAGS calculation, which continues to use the selected operation size.

Memory destination writes transfer only the selected number of bytes. Extension-store instructions instead write
the explicitly named W, L, or Q destination width. SP has only Q-sized writable forms. The result-extension rule
applies to Rn destinations.

\manualtablecaption{Integer Data Sizes}
\begin{manualtabular}{@{}>{\raggedright\arraybackslash}p{0.45in}>{\raggedright\arraybackslash}p{0.55in}>{\raggedright\arraybackslash}p{0.45in}>{\raggedright\arraybackslash}p{0.50in}>{\raggedright\arraybackslash}p{1.55in}@{}}
\toprule
\rowcolor{ManualHeaderFill}
\textbf{Code} & \textbf{Suffix} & \textbf{Bits} & \textbf{Bytes} & \textbf{Name}\\
\midrule
\texttt{B} & \texttt{@B_SUFFIX@} & @B_BITS@ & @B_BYTES@ & Byte\\
\texttt{W} & \texttt{@W_SUFFIX@} & @W_BITS@ & @W_BYTES@ & Word\\
\texttt{L} & \texttt{@L_SUFFIX@} & @L_BITS@ & @L_BYTES@ & Long\\
\texttt{Q} & \texttt{@Q_SUFFIX@} & @Q_BITS@ & @Q_BYTES@ & Quad\\
\bottomrule
\end{manualtabular}

\begin{manuallistedformatdiagram}{Integer Register Operand Subfields}
\manualformatrow{Q}{%
\manualformatfield{Q[63:0]}{64}
}
\manualformatrow{L}{%
\manualformatfield{unchanged / instruction-defined}{32}
\manualformatfield{L[31:0]}{32}
}
\manualformatrow{W}{%
\manualformatfield{unchanged / instruction-defined}{48}
\manualformatfield{W[15:0]}{16}
}
\manualformatrow{B}{%
\manualformatfield{unchanged / instruction-defined}{56}
\manualformatfield{B[7:0]}{8}
}
\end{manuallistedformatdiagram}


\subsection{Little-Endian Memory Organization}
Integer data, immediates, displacements, and absolute address payloads are little-endian. Multi-byte numeric
payloads are stored least significant byte first. The decoder consumes instruction-header fields byte by byte in
instruction-stream order.

For an N-byte integer value stored at byte address A, byte A contains bits 7..0, byte A+1 contains bits 15..8, and
so on.

\begin{manuallistedbyteorderdiagram}{Little-Endian Byte Order for a 64-Bit Value}{byte}{address}{increasing addresses}{least significant byte first}{most significant byte last}
\manualbytecell{$V[7..0]$}{A+0}
\manualbytecell{$V[15..8]$}{A+1}
\manualbytecell{$V[23..16]$}{A+2}
\manualbytecell{$V[31..24]$}{A+3}
\manualbytecell{$V[39..32]$}{A+4}
\manualbytecell{$V[47..40]$}{A+5}
\manualbytecell{$V[55..48]$}{A+6}
\manualbytecell{$V[63..56]$}{A+7}
\end{manuallistedbyteorderdiagram}

\begin{manuallistedbyteorderdiagram}{Instruction Start Bytes}{header byte}{}{instruction stream byte order}{}{}
\manualbytecell{\shortstack{class/length\\opcode high}}{PC+0}
\manualbytecell{\shortstack{opcode bits\\if present}}{PC+1}
\end{manuallistedbyteorderdiagram}


\subsection{Floating-Point Data Formats}
\manualtablecaption{Floating-Point Scalar Sizes}
\begin{manualtabular}{@{}>{\raggedright\arraybackslash}p{0.45in}>{\raggedright\arraybackslash}p{0.55in}>{\raggedright\arraybackslash}p{0.45in}>{\raggedright\arraybackslash}p{0.50in}>{\raggedright\arraybackslash}p{1.55in}@{}}
\toprule
\rowcolor{ManualHeaderFill}
\textbf{Code} & \textbf{Suffix} & \textbf{Bits} & \textbf{Bytes} & \textbf{Name}\\
\midrule
\texttt{S} & \texttt{@S_SUFFIX@} & @S_BITS@ & @S_BYTES@ & Single\\
\texttt{D} & \texttt{@D_SUFFIX@} & @D_BITS@ & @D_BYTES@ & Double\\
\bottomrule
\end{manualtabular}

Floating-point scalar formats are little-endian in memory as well. S is IEEE-754 binary32 with 24 bits of
significand precision, minimum normal exponent $-126$, and minimum subnormal quantum $2^{-149}$. D is IEEE-754
binary64 with 53 bits of significand precision, minimum normal exponent $-1022$, and minimum subnormal quantum
$2^{-1074}$.

\begin{manuallistedformatdiagram}{Floating-Point Register Encodings}[fig:floating-point-register-encodings]
\manualformatrowrange{D / Fn[63:0]}{63}{0}{%
\manualformatfield{s}{1}
\manualformatfield{exponent}{11}
\manualformatfield{fraction}{52}
}
\manualformatrowrange{S / Fn[63:0]}{63}{0}{%
\manualformatfield{zero}{32}
\manualformatfield{s}{1}
\manualformatfield{exponent}{8}
\manualformatfield{fraction}{23}
}
\end{manuallistedformatdiagram}
{\small
For a NaN, the most significant fraction bit is the quiet bit: bit 22 for S and bit 51 for D.
\texttt{s} denotes the sign bit. An S result occupies Fn bits 31..0 and writes zero to Fn bits 63..32.
\par}


The bit-level floating-point encoding, NaN policy, exception flags, and rounding behavior are defined by the
floating-point instruction and state descriptions; their memory byte order follows the same
least-significant-byte-first rule as integer data.
